\section{Experiments}


% sourced from DrivingWorld https://arxiv.org/pdf/2412.19505
\begin{table*}[t!]
\begin{minipage}[t]{\linewidth}
\caption{Multi-view Video Generation Performance.}
\label{tab:multiview}
\vspace{-10pt}
\centering
\begin{tabular}{lccccc}
\toprule
\multicolumn{1}{l}{\multirow{2}{*}{\textbf{Method}}} & \multirow{2}{*}{\textbf{Resolution}}& \multicolumn{3}{c}{\textbf{Metrics}}\\
& & FID$\downarrow$ & FVD$\downarrow$&Throughput$\uparrow$\\
\midrule
% % DriveSim~\cite{santana2016learningdrivingsimulator} & 7h & 80$\times$160    && & &\\
% % >>>>>> sourced from drivingwolrd
% DriveGAN~\cite{kim2021drivegancontrollablehighqualityneural} & \colorbox[gray]{0.9}{160h} & 256$\times$256    & GAN & 73.4& 502.3&-\\
% % ADriver-I~\cite{jia2023adriverigeneralworldmodel} & \colorbox[gray]{0.9}{300h} & 256$\times$512    &Diffusion& & &\\
% GenAD~\cite{yang2024genadgeneralizedpredictivemodel} & 2000h & 256$\times$448   &Diffusion& 15.4 & 184.0&-\\
% % GAIA-1~\cite{hu2023gaia1generativeworldmodel} & \colorbox[gray]{0.9}{4700h} & 288$\times$512   &Diffusion& - & - & - \\
% Vista~\cite{gao2024vistageneralizabledrivingworld} & 1740h & 576$\times$1024  &Diffusion& 6.9& 89.4&-\\
% % <<<<<< sourced from drivingwolrd
%  % DrivingGPT~\cite{chen2024drivinggptunifyingdrivingworld}& 128h & & AR& & &\\
%  % >>>>>> sourced from Driving into the Future
% WoVoGen~\cite{lu2023wovogen} & 5h & 256$\times$448    &Diffusion& 27.6 & 417.7&-\\
%   % <<<<<< sourced from Driving into the Future
% DrivingWorld*~\cite{hu2024drivingworldconstructingworldmodel} & 120h & 512$\times$1024   & AR & 16.4 & 174.4 & 0.42 \\
% DrivingWorld~\cite{hu2024drivingworldconstructingworldmodel} & \colorbox[gray]{0.9}{3456h} & 512$\times$1024   & AR & 7.4 & 90.9& 0.42\\
%   Ours (single-view) & 60h &  384$\times$672  &Next-scale AR& 4.5& &\\
%   \midrule
MagicDrive~\cite{gao2023magicdrive}& 224$\times$400&  16.2 & - & 1.76\\
X-Drive~\cite{xie2024x} &  224$\times$400 &  16.0 & - & 0.83\\
% Drive-WM~\cite{wang2023drivingfuturemultiviewvisual} & 192$\times$384 & 15.8&123 & -\\
DriveDreamer~\cite{wang2023drivedreamer} & 256$\times$448& 14.9 & 341& 0.37\\
BEVWorld~\cite{zhang2024bevworld} & - & 19.0 & 154 & -\\
Panacea~\cite{wen2024panacea} & 256$\times$512 & 17.0 & 139 & 0.67\\
DrivingDiffusion~\cite{li2024drivingdiffusion} & 512$\times$512 & 15.8 & 332 & -\\
\midrule
\ourmethod{} & 384$\times$672 & \textbf{10.5}& \textbf{91} & \textbf{1.96}\\
\bottomrule
\end{tabular}
\end{minipage}
% \begin{minipage}[t]{0.33\linewidth}
% \centering
% \caption{Front-View Performance}
% \label{tab:single}
% \vspace{-10pt}
% \begin{tabular}{cc}
% \toprule
% \textbf{Method}& \textbf{FID} $\downarrow$ \\
% \midrule
% % DriveSim~\cite{santana2016learningdrivingsimulator} & 7h & 80$\times$160    && & &\\
% % >>>>>> sourced from drivingwolrd
% DriveGAN~\cite{kim2021drivegancontrollablehighqualityneural} & 73.4\\
% % ADriver-I~\cite{jia2023adriverigeneralworldmodel} & \colorbox[gray]{0.9}{300h} & 256$\times$512    &Diffusion& & &\\
% GenAD~\cite{yang2024genadgeneralizedpredictivemodel}& 15.4 \\
% % GAIA-1~\cite{hu2023gaia1generativeworldmodel} & \colorbox[gray]{0.9}{4700h} & 288$\times$512   &Diffusion& - & - & - \\
% Vista~\cite{gao2024vistageneralizabledrivingworld} & 6.9\\
% WoVoGen~\cite{lu2023wovogen} &27.6 \\
%   % <<<<<< sourced from Driving into the Future
% DrivingWorld~\cite{hu2024drivingworldconstructingworldmodel} & 7.4 \\
% GEM~\cite{hassan2025gem} & 10.3\\
% GeoDrive~\cite{chen2025geodrive} & 4.1\\
% \midrule
% \ourmethod{} & \textbf{3.1}\\
% \bottomrule
% \end{tabular}
% \end{minipage}
\vspace{-10pt}
\end{table*}


\subsection{Experimental Setups}
\label{sec:setups}

\noindent\textbf{Datasets.} \ourmethod{} can scale up to heterogeneous training data with varying camera configurations, resolutions, or frame rates. We combine two public datasets (nuScenes~\cite{caesar2020nuscenes} and nuPlan~\cite{caesar2022nuplanclosedloopmlbasedplanning}), which are converted to a unified ScenarioNet~\cite{li2023scenarionet} format. NuScenes~\cite{caesar2020nuscenes} dataset contains $\sim$5 hours of driving data collected by 6 cameras with bounding box and map annotations, and scene descriptions from OmniDrive~\cite{wang2025omnidrive}. For nuPlan~\cite{caesar2022nuplanclosedloopmlbasedplanning} dataset, we filter the entire dataset based on complete sensor coverage and extract $\sim$55 hours of data from 8 cameras with auto-labeled bounding boxes and maps. We apply GPT-4o mini~\cite{team2024gpt} for scene descriptions. \ourmethod{} can be further scaled to larger-scale data, with details in the supplementary materials.
% Our in-house data include  $~$1M (overlapped) scenarios collected with 6 cameras in California and Michigan, with auto-labeled bounding boxes and without scene descriptions or HD-maps. However, 

\noindent\textbf{Evaluation Metrics.}  To evaluate synthetic image quality, we adopt Frechet Video Distance (FVD)~\cite{unterthiner2018towards} and Frechet Inception Distance (FID)~\cite{heusel2017gans}. We also evaluate fidelity to conditions through perception models using NDS for objects and mIoU for map. Unless otherwise specified, we conduct the evaluation on the nuScenes validation set~\cite{caesar2020nuscenes}.

\noindent\textbf{Implementation Details.} We build \ourmethod{} on two variants of the pretrained Infinity~\cite{han2024infinity} model with 130M and 2B parameters respectively. Our training pipeline includes three stages: 1) We start from training the model on low-resolution ($192\times336$) multi-view short clips 2) The model is further finetuned on short clips of high-resolution ($384\times672$) multi-view images. 3) Finally, we adopt the recurrent training in Sec.~\ref{sec:training} on long sequence of driving scenarios to learn the long-term temporal coherence. All the training is conducted on 32 NVIDIA A100 GPUs.

\subsection{Video Generation Performance}

\noindent\textbf{Results and Comparisons.} In Tab.~\ref{tab:multiview}, we compare the multi-view generation performance of \ourmethod{} with other existing world models. We do not include in-house data in our model training, which makes our training data scale less than most of other approaches. Results show that we attain FID and FVD much better than all the baselines, demonstrating better image quality and temporal coherence. Notably, \ourmethod{} delivers a throughput of \textbf{1.96 images/s}, much faster than diffusion baselines, demonstrating the efficiency of our autoregressive architecture. It is also worth mentioning that our model applies an image-based VAE encoder-decoder to minimize the inductive bias relevant to camera setups and ego-motions, although this design choice may hurt the FID and FVD metric. We also provide some qualitative visualizations in Fig. \ref{fig:teaser}. It can be observed that our model can generate under different conditions. Using dual-causal autoregression, \ourmethod{} not only matches or exceeds the performance of models trained on orders of magnitude more private data, but also operates at higher frame rates. 


\noindent\textbf{Object \& Map Conditions.} Our model takes optional object and map inputs as extra local geometric conditions to control the multi-view video generation. In addition to qualitative visualization in Fig.~\ref{fig:teaser}, we provide some additional quantitative results of object and map condition fidelity. We apply some state-of-the-art perception models~\cite{wang2023exploring,xie2023sparsefusion,liu2022bevfusion} to our synthetic images in nuScenes validation set. Both vision-based and multi-modal (combined with ground-truth point clouds) perception models are considered in the evaluation to comprehensively evaluate both the semantic and geometric realism. It is worth mentioning that StreamPETR takes long sequences as input, requiring  not only object realism in each single frame but also temporal coherence over long sequences. In Tab.~\ref{tab:object_condition} and \ref{tab:map_condition}, our synthetic images show great fidelity to both objects and maps notably better than baselines. For map condition, the projection of coordinates to camera space is not perfect due to the lack of height information in the map annotation, but \ourmethod{} still shows robust performance. For object condition, our synthetic images can reach a performance similar to real images.

\begin{table}[tb]
        \caption{Fidelity to Object Condition}
        \label{tab:object_condition}
        \vspace{-10pt}
        \begin{subtable}[t]{0.48\linewidth}
            \centering
            \subcaption{Camera (StreamPETR\cite{wang2023exploring})}
            \label{tab:object_camera} 
            \resizebox{\linewidth}{!}{
            \begin{tabular}{c|c}
                \toprule
                \textbf{Method} & \textbf{NDS$\uparrow$}  \\
                \midrule
                \textcolor{gray}{Oracle} & \textcolor{gray}{46.9} \\
                \midrule
                Panacea~\cite{wen2024panacea} & 32.1 \textcolor{red}{(68\%)}\\
                \ourmethod{} & \textbf{41.9} \textcolor{darkgreen}{(89\%)}\\
                \bottomrule
            \end{tabular}
        }
        \end{subtable}
        \begin{subtable}[t]{0.48\linewidth}
            \centering
            \subcaption{Multisensor (SparseFusion\cite{xie2023sparsefusion})}
            \label{tab:object_fusion} 
            \resizebox{\linewidth}{!}{
            \begin{tabular}{c|c}
                \toprule
                \textbf{Method} & \textbf{NDS$\uparrow$}  \\
                \midrule
                \textcolor{gray}{Oracle} & \textcolor{gray}{72.8} \\
                \midrule
                X-Drive~\cite{xie2024x} & 69.6 \textcolor{red}{(95\%)}\\
                \ourmethod{} & \textbf{72.0} \textcolor{darkgreen}{(99\%)}\\
                \bottomrule
            \end{tabular}
            }
        \end{subtable}
\end{table}

\begin{table}[tb]
        \vspace{-10pt}
        \caption{Fidelity to Map Condition}
        \label{tab:map_condition}
        \vspace{-10pt}
        \begin{subtable}[t]{0.48\linewidth}
            \centering
            \subcaption{Camera (BEVFusion-C\cite{liu2022bevfusion})}
            \label{tab:map_camera} 
            \resizebox{\linewidth}{!}{
            \begin{tabular}{c|c}
                \toprule
                \textbf{Method} & \textbf{mIoU$\uparrow$}  \\
                \midrule
                \textcolor{gray}{Oracle} & \textcolor{gray}{57.1} \\
                \midrule
                MagicDrive~\cite{gao2023magicdrive} & 28.9 \textcolor{red}{(51\%)}\\
                \ourmethod{} & \textbf{34.9} \textcolor{darkgreen}{(61\%)}\\
                \bottomrule
            \end{tabular}
        }
        \end{subtable}
        \begin{subtable}[t]{0.48\linewidth}
            \centering
            \subcaption{Multisensor (BEVFusion\cite{liu2022bevfusion})}
            \label{tab:map_fusion} 
            \resizebox{\linewidth}{!}{
            \begin{tabular}{c|c}
                \toprule
                \textbf{Method} & \textbf{mIoU$\uparrow$}  \\
                \midrule
                \textcolor{gray}{Oracle} & \textcolor{gray}{63.0} \\
                \midrule
                MagicDrive~\cite{gao2023magicdrive} & 47.0 \textcolor{red}{(75\%)}\\
                \ourmethod{} & \textbf{49.9} \textcolor{darkgreen}{(79\%)}\\
                \bottomrule
            \end{tabular}
            }
        \end{subtable}
    \vspace{-10pt}
\end{table}

\begin{table}[tb!]
    \centering
    \caption{Motion Cues in Generated Videos}
    \label{tab:motion}
    \vspace{-10pt}
    \resizebox{0.9\linewidth}{!}{
    \begin{tabular}{c|ccc}
        \toprule
       \multirow{2}{*}{\textbf{Method}}  & \multicolumn{3}{c}{\textbf{Motion Planning} ($L_2\downarrow$)}  \\
         & $1s$ & $2s$ & $3s$\\
        \midrule
        \textcolor{gray}{Oracle} & \textcolor{gray}{0.41} & \textcolor{gray}{0.70} & \textcolor{gray}{1.05}\\
        \midrule
        \ourmethod{} & \textbf{0.42} \textcolor{darkgreen}{(98\%)} & \textbf{0.73} \textcolor{darkgreen}{(96\%)} & \textbf{1.09} \textcolor{darkgreen}{(96\%)}\\
        \bottomrule
    \end{tabular}
    }
\end{table}



\begin{table*}[htb]
    \centering
    \caption{Novel View Synthesis Performance with Camera Shifts.}
    \label{tab:nvs}
    \vspace{-10pt}
    \begin{tabular}{l|cc|cc|cc}
        \toprule
        \multirow{3}{*}{Method} & \multicolumn{2}{c|}{Shift $1m$} & 
        \multicolumn{2}{c|}{Shift $2m$} & \multicolumn{2}{c}{Shift $4m$}\\
        \cmidrule(lr){2-3} \cmidrule(lr){4-5} \cmidrule(lr){6-7}
        & FID$\downarrow$ & FVD$\downarrow$ & FID$\downarrow$ & FVD$\downarrow$ &FID$\downarrow$ &FVD$\downarrow$\\
        \midrule
        % PVG~\cite{chen2023periodic}      &  48.15 & 246.74  & 60.44 & 356.23 & 84.50 & 501.16  \\
        % EmerNeRF~\cite{yang2023emernerf}   &  37.57  &  171.47  & 52.03  & 294.55 & 76.11 &  497.85 \\ 
        StreetGaussian~\cite{yan2024street} & 32.12  & 153.45  &  43.24 & 256.91 & 67.44 & 429.98 \\
        OmniRe~\cite{chen2024omnire}   & 31.48  & 152.01  & 43.31  & 254.52 & 67.36 & 428.20 \\
        % \rowcolor{gray!10}w/. Ours & \textbf{81.69} & \textbf{41.62}  & \textbf{36.50} & \textbf{41.22} \\ 
        % \midrule
        FreeVS~\cite{wang2024freevs}   &  51.26  &  431.99  &  62.04  &  497.37  &  77.14  &  556.14 \\
        % ViewCrafter~\cite{yu2024viewcrafter}   &  14.84  & - &  17.15  & -   &  18.98  & - \\
        \midrule
        \ourmethod{} & \textbf{14.11}  & \textbf{117.20} & \textbf{15.58} & \textbf{122.60} & \textbf{17.48} & \textbf{143.48}\\
        \bottomrule
    \end{tabular}
\end{table*}

\noindent\textbf{Motion Cues in Synthetic Videos.} In addition to visual quality and condition fidelity, we further evaluate whether our synthetic videos support plausible driving decisions by applying an end-to-end planner (VAD~\cite{jiang2023vad}) pretrained on real scenes to our generated videos. Tab.~\ref{tab:motion} shows that the pretrained planner derives actions from synthetic videos consistent with real scenes. This reflects not only visual realism but also physical and motion plausibility in our generated videos. 

\noindent\textbf{Zero-shot to Unseen Camera Configurations.} The ray-level relative position embedding allows \ourmethod{} to generalize to unseen camera configurations in a zero-shot manner. An example of sensor setups different from nuScenes and nuPlan is visualized in Fig.~\ref{fig:waymo} of the appendix.




\subsection{Novel View Synthesis}
\label{sec:nvs}
The relative position embedding in camera ray space allows our model to synthesize images for novel views. We follow \cite{wang2024freevs} to evaluate the performance of \ourmethod{} for novel view synthesis task on nuScenes validation set. Different shifts $\{1m,2m,4m\}$ are applied to the camera viewpoints. We measure the FID and FVD metrics between the synthetic shifted multi-view videos and the original real videos. In Tab.~\ref{tab:nvs}, our FID and FVD metrics only drop slightly after viewpoint shift. Compared to baselines, \ourmethod{} exhibits a significant performance gain even if it does not include any 3D inductive bias such as radiance field or point clouds. These results demonstrate the great capability of \ourmethod{} in camera controllability and novel view synthesis in a geometry-free data-driven manner.

\begin{table}[tb]
    \centering
    \vspace{-10pt}
    \caption{Ablation Study on Camera Ray.}
    \label{tab:position}
    \vspace{-10pt}
    \begin{tabular}{cc|cc}
    \toprule
        \textbf{S.-T. Module} & \textbf{Pos. Emb.} & \textbf{FID$\downarrow$} & \textbf{FVD$\downarrow$}  \\
        \midrule
         \XSolidBrush & - &  18.7 & 214 \\
         \Checkmark & Absolute & \textbf{17.2} & 138 \\
         \Checkmark & Relative & \textbf{17.2} & \textbf{124} \\
    \bottomrule
    \end{tabular}
\end{table}

\begin{figure}[tb!]
    \centering
    \includegraphics[width=\linewidth]{figs/ablation_scale.pdf}
    \vspace{-20pt}
    \caption{Ablation Study on Scale Causality. Conditioning on all scales in history hurts the modeling of dynamics, while conditioning only on same scale is insufficient for temporal coherence.}
    \label{fig:casuality}
    \vspace{-5pt}
\end{figure}


\begin{table}[tb!]
    \centering
    \caption{Ablation Study on Scale Causality and Model Size.}
    \label{tab:casuality}
    \vspace{-10pt}
    \begin{tabular}{cccc}
        \toprule
         \textbf{Past-frame Condition} & \textbf{Size} & \textbf{FID$\downarrow$} & \textbf{FVD$\downarrow$}  \\
        \midrule
         same scale & 130M & 20.7 & 151 \\
         all scales & 130M &  60.4 & 454 \\
         \midrule
         prefix scales & 130M & 17.2 & 124 \\
         \midrule
         prefix scales & 2B & \textbf{10.5} & \textbf{91}\\
        \bottomrule
    \end{tabular}
    \vspace{-15pt}
\end{table}

\begin{figure}[tb]
    \centering
    \includegraphics[width=\linewidth]{figs/model_scale.pdf}
    \vspace{-20pt}
    \caption{Ablation Study on Model Size. Large scale model can bring significantly better visual quality.}
    \label{fig:model_scale}
    \vspace{-10pt}
\end{figure}

\subsection{Ablation Study}
\noindent\textbf{Scale Causality.}
\ourmethod{} is built up on scale and time dual-causality. The temporal causality is 
inevitable if we would like to extend to long video autoregression. For scale causality (Eq.~\ref{eq:formulation}), different from low-scale condition in previous frames, we consider two alternatives conditions: 1) all scales or 2) same scale features in previous frames. Results are shown in Tab.~\ref{tab:casuality} and Fig.~\ref{fig:casuality}. Higher scales features of previous time steps cause future frames simply copying the details without simulating the dynamics, and same scale condition is far from enough for stable spatio-temporal modeling.

\noindent\textbf{Relative Ray Embedding.}
One of our key contributions lies in the relative position encoding in the camera ray space, which is critical for the spatio-temporal consistency. To justify our design, we consider several alternatives in Tab.~\ref{tab:position}. We find that this relative position embedding can perform better than absolute Pl\"ucker ray embedding, let alone none spatio-temporal reliance.

\noindent\textbf{Recurrent Training.} The recurrent training stage is greatly helpful for temporal coherence in video generation. It improves FVD from 100 to \textbf{91} within hundreds of iterations.

\noindent\textbf{Model Size.} Our minimal reliance on inductive bias provides a data-driven framework. Large model can exhibit significantly superior performance. We compare two variants of our model with 130M and 2.8B parameters. In Fig.~\ref{fig:model_scale} and Tab.~\ref{tab:casuality}, with the same training data, larger model enables much better image quality.

